% chap-literature.tex  –  Chapter 2: Literature Review and Theoretical Framework

\chapter{Literature Review and Theoretical Framework}
\label{chap:literature}

%% ----------------------------------------------------------------
\section{Historical Automation Risk Research}
\label{sec:historical-automation}

The modern study of automation risk begins with \citet{Frey_Osborne_2017}'s landmark paper. Working with occupational data from the U.S.\ Department of Labor's O*NET database, they developed a probabilistic framework for estimating how susceptible individual occupations were to computerization. Their approach was straightforward: they first had machine learning researchers at Oxford hand-label 70 occupations as either fully automatable or not automatable based on professional judgment about the technological frontier in 2013. They then identified nine O*NET variables corresponding to three engineering bottlenecks (perception and manipulation (e.g., finger dexterity, manual dexterity), creative intelligence (e.g., originality, fine arts), and social intelligence (e.g., social perceptiveness, negotiation, persuasion)) and trained a Gaussian Process Classifier to predict automation probability for all 702 occupations in their dataset. The result was a probability between 0 and 1 for each occupation, representing the likelihood that computerization would automate that job within the next two decades.

Their central result: 47\% of total U.S.\ employment fell into the high-risk category, with automation probabilities exceeding 0.70. The occupations at greatest risk were dominated by routine cognitive and manual tasks: telemarketers (0.99), tax preparers (0.99), and hand sewers (0.99). At the other extreme, occupations requiring intensive creative or social intelligence (recreational therapists (0.003), mental health counselors (0.004), and choreographers (0.004)) were rated as essentially immune. The three bottlenecks formed the theoretical boundary between human work and machine capability.

\citet{Arntz_Gregory_Zierahn_2016} challenged Frey and Osborne's methodology in an OECD working paper. They argued that classifying entire occupations as automatable or not was too coarse, because workers within the same occupation perform different mixes of tasks. Using the PIAAC (Programme for the International Assessment of Adult Competencies) survey, they analyzed task-level variation within occupations and found that only 9\% of jobs across 21 OECD countries were at high risk of automation when this variation was considered, dramatically less than Frey and Osborne's 47\%. This critique is directly relevant to my project: it supports treating automation exposure as a continuous spectrum rather than a binary outcome, which is exactly how the MAEI's feature-level multiplier approach operates.

\citet{Nedelkoska_Quintini_2018} extended the \citet{Arntz_Gregory_Zierahn_2016} methodology to 32 countries using PIAAC data and found that 14\% of jobs across OECD nations were highly automatable, with another 32\% facing significant changes in task composition. Their cross-national analysis revealed substantial variation in automation vulnerability driven by differences in industrial composition and educational attainment. Countries with larger manufacturing sectors and lower average educational levels showed higher shares of automatable employment. This geographic variation confirms that any single-country index like the MAEI captures only one slice of a globally heterogeneous phenomenon, but the U.S.\ case remains the most data-rich and therefore the most suitable for methodological development.

The transition from the historical automation literature to the modern AI exposure literature reflects a fundamental shift in the nature of the technological threat. Where \citet{Frey_Osborne_2017} asked whether a robot could perform a job's physical tasks, the relevant question in 2026 is whether a language model can perform a job's cognitive tasks. This reframing does not invalidate the earlier research; physical automation continues advancing in warehouse logistics and food service, but it demands a complementary analytical framework that captures the new cognitive dimension. The MAEI provides this complementary framework by preserving the historical structure while overlaying a modern capability assessment calibrated to the LLM era.

%% ----------------------------------------------------------------
\section{The Task-Based Economic Framework}
\label{sec:task-based}

\citet{Acemoglu_Restrepo_2018, Acemoglu_Restrepo_2019} provided the economic theory that justifies disaggregating occupations into tasks. Their task-based framework argues that automation does not destroy ``jobs'' wholesale; instead, it substitutes capital for labor at the level of individual tasks within occupations. When a firm automates one task (say, data entry), the worker's other tasks (client communication, quality review) may become more valuable. Their model also shows that automation triggers two opposing forces: a displacement effect (machines replace workers in certain tasks, reducing labor demand) and a reinstatement effect (new tasks emerge that require human input, increasing labor demand). The net impact depends on the relative strength of these forces.

\citeauthor{Acemoglu_Restrepo_2018}'s framework informed my design choices. I relied on O*NET's granular task-level data rather than treating occupations as monolithic units. My capability multipliers target specific features (e.g., Written Expression, Programming, Mathematical Reasoning) rather than entire occupations, reflecting the task-based insight that AI automates specific capabilities, not entire jobs. The distinction between task exposure and job elimination is also why I insist that the MAEI measures capability overlap, not employment loss.

%% ----------------------------------------------------------------
\section{Generative AI as General-Purpose Technology}
\label{sec:gpt-technology}

\citet{Eloundou_Manning_2023}, working at OpenAI, reframed the automation question for the LLM era. They defined large language models as ``General-Purpose Technologies'' (GPTs), a term from innovation economics meaning technologies that affect virtually every sector of the economy, not just specific industries. By combining human and GPT-4 assessments of which O*NET tasks could be performed or accelerated by LLMs, they estimated that approximately 80\% of the U.S.\ workforce could have at least 10\% of their work tasks affected by LLMs. For about 19\% of workers, the impact could extend to over 50\% of their tasks.

That finding directly influenced my project. If 80\% of the workforce faces some degree of LLM exposure, then any index limited to the old F\&O model (which concentrated risk on routine manual work) would fundamentally undercount modern exposure. This is exactly what I observed when running my baseline model: even with aggressive cognitive multipliers, the 2013-trained model resisted assigning high exposure to white-collar occupations because it had learned that cognitive skills were ``protective.'' \citet{Eloundou_Manning_2023}'s 80\% figure motivated my Broad Exposure Uplift formula, which mathematically corrects for this structural bias.

\citet{Noy_Zhang_2023} provided experimental evidence complementing Eloundou's theoretical estimates. In a randomized controlled trial with 453 college-educated professionals, they found that access to ChatGPT reduced average task completion time by 37\% and improved output quality as rated by blind evaluators. The most notable finding was an inequality-compressing effect: below-average writers benefited far more than skilled writers, suggesting that AI acts as a productivity equalizer in cognitive work. I used these experimental findings to justify the specific multiplier magnitudes for language-related O*NET features (Written Expression at $2.80\times$, Written Comprehension at $2.70\times$, Reading Comprehension at $2.60\times$).

%% ----------------------------------------------------------------
\section{Modern AI Exposure Measurement}
\label{sec:modern-measurement}

\citet{Felten_Raj_Seamans_2021} constructed the AI Occupational Exposure (AIOE) index using a methodologically distinct approach. They crowdsourced expert linkages between specific AI applications (image recognition, translation, speech recognition, and seven others) and O*NET occupational abilities. By weighting each occupation's abilities by the corresponding AI application's capability, they produced a continuous exposure score for 782 occupations. Their most notable finding was that high-wage, highly-educated white-collar occupations showed the greatest AI exposure, a direct inversion of \citet{Frey_Osborne_2017}'s 2013 profile. I used the AIOE dataset as my external validation benchmark specifically because its methodology is completely independent of mine: they used crowdsourced human judgments while I used scenario-based capability multipliers.

\citet{Webb_2020} took yet another approach, using natural language processing to match the text of AI-related patents to the text of occupational task descriptions. By measuring the semantic overlap between what AI patents claim to do and what workers actually do, he showed that AI technologies disproportionately target high-skill, high-wage cognitive work. This validated my decision to use NLP (specifically BERT sentence embeddings) to extract latent semantic features from O*NET task descriptions, capturing dimensions of occupational content that numeric scales alone cannot represent.

Most recently, \citet{Anthropic_2026} published research measuring ``observed exposure'' by combining real-world Claude usage data with O*NET task descriptions and \citet{Eloundou_Manning_2023}'s feasibility ratings. They found that Computer Programmers showed the highest coverage at 75\%, that AI-exposed workers earn 47\% more on average than unexposed workers, and, critically, that there was ``no systematic unemployment increase'' for workers in highly exposed occupations since late 2022. While their approach uses proprietary usage data that cannot be replicated in academic settings, their findings independently corroborate the MAEI's core thesis: cognitive, high-wage occupations face the greatest AI capability overlap, but this overlap has not yet translated into measurable job displacement.

The progression from \citet{Frey_Osborne_2017} through \citet{Felten_Raj_Seamans_2021}, \citet{Eloundou_Manning_2023}, and \citet{Anthropic_2026} reveals a clear methodological arc in the field. Each successive study has broadened the definition of AI exposure, shifted the locus of vulnerability from manual to cognitive work, and employed increasingly sophisticated measurement approaches. Yet no single study has combined all the methodological elements needed for a comprehensive, reproducible, and externally validated index built entirely from public data. The MAEI draws structural inspiration from \citet{Frey_Osborne_2017}, validation methodology from Felten, breadth estimation from Eloundou, and empirical grounding from Anthropic's observed usage patterns, integrating these strands into a single reproducible pipeline.

%% ----------------------------------------------------------------
\section{Creativity and Human-Centric Value}
\label{sec:creativity}

A persistent question in automation research is whether the creative intelligence bottleneck identified by \citet{Frey_Osborne_2017} still holds in the age of generative AI. \citet{Doshi_Hauser_2024} conducted a large-scale experiment examining AI's effect on creative work and found a mixed answer: generative AI enhances individual creativity (helping people produce more novel ideas) but reduces the collective novelty of creative output (everyone converges on similar AI-assisted solutions). This suggests that while AI can mimic creative output at the individual task level, the uniquely human capacity for sustained, divergent originality remains distinctive.

My model's feature importances support this interpretation. The top two most important features in the baseline Stacked Ensemble are Fluency of Ideas (importance = 0.2177) and Originality (importance = 0.1919), which together account for over 40\% of the model's predictive power. These are precisely \citet{Frey_Osborne_2017}'s creative intelligence bottleneck variables. Occupations scoring highly on these features receive the lowest MAEI scores, suggesting that deep creative capability remains the strongest structural defense against AI exposure.

%% ----------------------------------------------------------------
\section{Research Gap}
\label{sec:research-gap}

To summarize: the existing literature on automation and AI exposure either (a) reflects pre-LLM technology and cannot capture the cognitive shift \citep{Frey_Osborne_2017}, (b) relies on proprietary usage data that cannot be replicated \citep{Anthropic_2026}, (c) measures exposure through crowdsourced judgments without scenario modeling or uncertainty quantification \citep{Felten_Raj_Seamans_2021}, or (d) provides theoretical economic frameworks without constructing occupation-level empirical indices \citep{Acemoglu_Restrepo_2018}. No published work combines an updated Frey and Osborne structural framework with BERT-based NLP task embeddings, literature-calibrated capability multipliers, Monte Carlo uncertainty bounds, and external validation against an independent benchmark, all built entirely from publicly available data. The MAEI addresses this gap.
